%% Chapter 12 — Verification Ladder and Falsifiers

\chapter{Verification Ladder and Falsifiers}
\label{chap:verification}

\epigraph{%
  ``For every complex problem there is an answer that is clear, simple, and wrong.''
}{--- H.L.\ Mencken; the verification ladder prevents premature closure}

\noindent Any formal system that cannot be falsified is not science — it is creed. This
chapter provides the two instruments that keep the \emph{Periodic Table of Reason} honest:
a ten-rung (plus one extension) verification ladder that any evaluator can climb
independently, and nine falsifiers that would, if realised, require the thesis to be
retracted or substantially revised. The ladder specifies \emph{what to run}; the
falsifiers specify \emph{what would constitute refutation}.

%%--------------------------------------------------------------------
\section{Verification Ladder}
\label{sec:verification-ladder}
%%--------------------------------------------------------------------

The verification ladder is ordered from least to most computationally demanding. An
evaluator who completes all rungs has independently reproduced every empirical claim in
the thesis.

\medskip
\textbf{1.} Clone the repository and confirm the monorepo structure matches
Section~\ref{sec:repository-structure}: \texttt{wasm4pm/} (Rust/WASM core),
\texttt{packages/} (11 TypeScript packages), \texttt{apps/wasm4pm/} (primary CLI),
\texttt{crates/} (secondary Rust crates including \texttt{wasm4pm-cognition},
\texttt{prolog8}, \texttt{miniml-core}, \texttt{ocpq}). Confirm that
\texttt{Cargo.toml} lists the workspace members and \texttt{pnpm-workspace.yaml}
lists the package paths.

\medskip
\textbf{2.} Build the WASM core: from \texttt{crates/wasm4pm-cognition/}, run
\begin{center}
\texttt{wasm-pack build --target nodejs --out-dir pkg -- --features wasm}
\end{center}
and confirm that \texttt{pkg/wasm4pm\_cognition\_bg.wasm} is produced without errors.
This rung verifies that the Rust source compiles to a valid WASM binary and that the
\texttt{wasm\_bindgen} exports (\texttt{cognition\_run}, \texttt{cognition\_verify},
\texttt{system\_build}) are present in the generated \texttt{.d.ts} stubs.

\medskip
\textbf{3.} Run the Rust unit test suite: from the repository root, run
\begin{center}
\texttt{cargo test 2>\&1 | grep -E "test result|FAILED|passed"}
\end{center}
and confirm 483 tests pass with 0 failures. The MYCIN CF boundary tests
(\texttt{test\_cf\_threshold\_boundary\_above}, \texttt{test\_cf\_threshold\_boundary\_below}
in \texttt{production\_rules.rs}) and the Prolog grandparent test
(\texttt{strips\_soar\_cbr\_invariants.rs}) must appear in the output. Any failure
triggers Andon (Section~\ref{sec:falsifiers}).

\medskip
\textbf{4.} Run the TypeScript test suite: from the repository root, run
\begin{center}
\texttt{pnpm build \&\& pnpm test}
\end{center}
and confirm all packages report passing. Specifically verify that
\texttt{packages/cognition/src/bvc.ts} (\texttt{computeBVC}, \texttt{isBVCCertified})
and \texttt{packages/contracts/src/admission.ts} (\texttt{BVCAdmissionResult},
\texttt{admitBVC}) are covered by passing tests.

\medskip
\textbf{5.} Run the Criterion latency benchmarks: from
\texttt{crates/wasm4pm-cognition/}, run
\begin{center}
\texttt{cargo bench --bench breed\_latency 2>\&1 | grep -E "time:|FAILED"}
\end{center}
and confirm that all 13 breeds report median latency $\leq 10\,\mu\mathrm{s}$.
Compare against \texttt{docs/benchmarks/breed-latency-2026-06-10.md}. Any breed
exceeding $10\,\mu\mathrm{s}$ is a candidate for Falsifier F6.

\medskip
\textbf{6.} Verify the BVC for each breed individually: invoke \texttt{wpm cognition run}
(or \texttt{wpm cr} alias) for each of the 13 breeds with a canonical test input.
Confirm that each response contains \texttt{status: "ok"}, a non-empty
\texttt{output\_hash}, a non-empty \texttt{input\_hash}, and that
\texttt{admitBVC(result)} returns \texttt{certified: true}. Any breed returning
\texttt{status !== "ok"} or \texttt{certified: false} is a candidate for
Falsifier F9.

\medskip
\textbf{7.} Verify OCEL conformance for one breed end-to-end: run \texttt{wpm cognition run}
for the \texttt{mycin} breed, extract the OTel spans, convert to OCEL format using
the \texttt{@wasm4pm/observability} span-to-ocel utility, and replay the resulting
log against the \MYCIN{} Petri net model using native Rust DFA replay via \texttt{validate\_ocel\_alignment()} (\texttt{ocel.rs}) against the declared lifecycle OCPN---deterministic, reproducible, no external process-mining framework. Confirm fitness $\phi = 1.0$.
Repeat for \texttt{prolog} as a second data point.

\medskip
\textbf{8.} Verify receipt chain integrity: run the factory-agent chain
(\texttt{examples/cognition/chains/factory-agent/chain.sh}) and confirm that all 13
stage receipts are produced and that
$r_{k+1}.\texttt{input\_hash} = r_k.\texttt{output\_hash}$ for $k = 0, \ldots, 11$.
The final receipt (\texttt{stages/12-eliza/result.json}) must contain a non-empty
\texttt{output\_hash} that is the BLAKE3 hash of the ELIZA stage's
\texttt{ocel\_log} concatenated with its \texttt{result}.

\medskip
\textbf{9.} Verify determinism for each breed: run the double-run bit-exact harness
\begin{center}
\texttt{cargo test breed\_determinism 2>\&1 | grep -E "test result|FAILED"}
\end{center}
and confirm all 13 breeds produce identical byte sequences on both runs with the
same input seed. Any non-determinism is Falsifier F3.

\medskip
\textbf{10.} Verify the \texttt{admitBVC} gate prevents non-certified breeds from entering
a pipeline: construct a synthetic \texttt{BVCAdmissionResult} with
\texttt{certified: false} and confirm that \texttt{admitBVC()} returns an error
response. Then construct one with \texttt{certified: true} and confirm it returns
an admit response. This rung verifies that the Fourth Constraint is enforced at
the API boundary, not merely documented.

\medskip
\textbf{10b.} Run all three breed chains via the master runner:
\begin{center}
\texttt{bash examples/cognition/chains/run-all-chains.sh}
\end{center}
Verify that all 26 stages across the three chains
(\texttt{socratic-diagnosis}: 7 breeds; \texttt{scientific-discovery}: 6 breeds;
\texttt{factory-agent}: 13 breeds) complete with exit code 0 and that every
stage receipt is BVC-certified. The expected summary line is:
\begin{center}
\texttt{PASSED: 3 / 3 chains}
\end{center}
Any chain failure is Andon; investigate the failing stage's receipt before
proceeding.

%%--------------------------------------------------------------------
\section{Falsifiers}
\label{sec:falsifiers}
%%--------------------------------------------------------------------

The following conditions, if demonstrated, would falsify one or more of the seven main
theorems. Falsification of a main theorem requires either retraction, revision, or a proof
that the falsifying instance lies outside the scope of the theorem's hypotheses.

\medskip
\textbf{F1 — Oracle forgery.} A test harness that does not implement the correct
breed algorithm nonetheless passes a Tier-2 oracle check. This would falsify
Theorem~\ref{thm:main-oracle} and would require reconstructing the oracle hierarchy
to identify the exploited weakness. A concrete demonstration requires constructing a
function $f \colon \mathcal{X} \to \mathcal{Y}$ that (a) does not implement \MYCIN{}
certainty-factor algebra, and (b) satisfies $\CF(A \wedge B) = \CF(A) \cdot \CF(B)$
for all inputs in the test suite. If such a function exists and is not the correct
algorithm, the oracle is not Tier-2.

\medskip
\textbf{F2 — OCEL fitness gap.} A breed's execution traces produce an OCEL event log
whose fitness under native Rust DFA replay via \texttt{validate\_ocel\_alignment()} (\texttt{ocel.rs}) against the declared lifecycle OCPN---deterministic, reproducible, no external process-mining framework, is $\phi < 1.0$. This would falsify
Theorem~\ref{thm:main-ocel}(iii) for that breed and would require revising either the
breed's process model $M_i$ or the OTel span-to-event mapping. The gap would also
reduce $C_{\mathrm{ocel}}(\Breed_i)$ below 1 and thus $V(\Breed_i)$ below 1.

\medskip
\textbf{F3 — Non-determinism.} Two invocations of the same breed with the same input
seed produce different byte-level outputs. This would falsify
Theorem~\ref{thm:main-bvc} ($C_{\mathrm{determ}} = 1$) for that breed and would
require identifying the non-determinism source: hash-map iteration order, floating-point
rounding differences across platforms, or an unseeded PRNG.

\medskip
\textbf{F4 — Receipt hash mismatch.} The \texttt{output\_hash} emitted by
\texttt{cognition\_run()} does not equal
$\blake3(\texttt{ocel\_log} \Vert \texttt{result})$ when recomputed independently.
This would falsify Theorem~\ref{thm:main-bvc} ($C_{\mathrm{receipt}} = 1$) and the
non-repudiation property of Theorem~\ref{thm:main-civ}. It would require a bug in
the BLAKE3 construction at the WASM boundary
(\texttt{crates/wasm4pm-cognition/src/wasm.rs}).

\medskip
\textbf{F5 — Epistemic dependence.} A functor $F \colon \prod_{j \neq i} \Breed_j
\to \Breed_i$ is exhibited in the \textbf{Breed} category that commutes with the
evidence functor, demonstrating that breed $\Breed_i$ is derivable from the remaining
twelve. This would falsify Theorem~\ref{thm:main-epistemic} for that breed and would
require either removing the breed from the table as redundant or strengthening the
independence proof. The most plausible candidate for approximate dependence is \ELIZA{}
(as a degenerate \CBR{} with empty case base), but the evidence functors differ in the
dimensions they populate in the epistemological space $\mathcal{V}$.

\medskip
\textbf{F6 — Latency regression above threshold.} A breed exceeds $10\,\mu\mathrm{s}$
median latency on representative hardware under Criterion benchmarking. This would
falsify Theorem~\ref{thm:main-speed} for that breed and would require either
performance optimisation or a revision of the phase-transition threshold argument.
Given that the current worst case is \PROLOG{} at $7.932\,\mu\mathrm{s}$, a regression
would require approximately a 26\% slowdown.

\medskip
\textbf{F7 — Regulatory inadmissibility finding.} A court, regulator, or standards body
issues a ruling that BLAKE3-based receipt chains are inadmissible as documentary evidence
in the relevant jurisdiction. This would falsify the regulatory-admissibility component
of Theorem~\ref{thm:main-civ} and would require replacing BLAKE3 with an
admissibility-certified hash function (e.g.\ SHA-3/256, which is currently FIPS 202
certified).

\medskip
\textbf{F8 — Fourth constraint sufficiency counter-example.} A pipeline satisfying
$C_0 \wedge C_1 \wedge C_2 \wedge C_3$ produces an output that is epistemically
untrustworthy in a way not detectable by the BVC. This would require constructing a
breed that achieves $V(\Breed_i) = 1$ by passing all four component tests while
concealing a systematic bias or error in its reasoning not exercised by the oracle suite.
A successful construction would require either expanding the oracle hierarchy to Tier-3
or adding a fifth constraint.

\medskip
\textbf{F9 — BVC formula gap.} The function \texttt{computeBVC()} in
\texttt{packages/cognition/src/bvc.ts} returns a value strictly less than 1.0 for a
breed invocation that is independently verified to satisfy all four components
($C_{\mathrm{test}} = C_{\mathrm{ocel}} = C_{\mathrm{receipt}} = C_{\mathrm{determ}} = 1$).
This would falsify Theorem~\ref{thm:main-bvc} not at the mathematical level but at the
implementation level, indicating a bug in the composite scoring formula:
\begin{center}
\texttt{score = (c\_test + c\_ocel + c\_receipt + c\_determ) / 4}
\end{center}
Specifically, this falsifier would be instantiated if floating-point summation error
caused \texttt{score} to underflow below 1.0 for four unit inputs, or if any of the
four component fields were incorrectly extracted from the \texttt{ContractResult}
returned by \texttt{cognition\_run()}.

%%--------------------------------------------------------------------
\section{Relationship Between Ladder and Falsifiers}
\label{sec:ladder-falsifiers-relationship}
%%--------------------------------------------------------------------

Each rung of the verification ladder is designed to exercise one or more falsifiers.
The correspondence is as follows:

\medskip
\begin{center}
\begin{tabular}{cll}
\toprule
\textbf{Rung} & \textbf{Primary Falsifier(s) Exercised} & \textbf{Theorem at Risk} \\
\midrule
1   & (structural, no direct falsifier)      & Background \\
2   & (compilation; no direct falsifier)     & Background \\
3   & F1 (oracle forgery)                    & Thm.~\ref{thm:main-oracle} \\
4   & F9 (BVC formula gap)                   & Thm.~\ref{thm:main-bvc} \\
5   & F6 (latency regression)                & Thm.~\ref{thm:main-speed} \\
6   & F4 (receipt hash mismatch), F9         & Thm.~\ref{thm:main-bvc}, \ref{thm:main-civ} \\
7   & F2 (OCEL fitness gap)                  & Thm.~\ref{thm:main-ocel} \\
8   & F4 (receipt hash mismatch)             & Thm.~\ref{thm:main-civ} \\
9   & F3 (non-determinism)                   & Thm.~\ref{thm:main-bvc} \\
10  & F8 (fourth constraint gap)             & Thm.~\ref{thm:main-fourth} \\
10b & F2, F4, F9 (chain-level integration)   & Thm.~\ref{thm:main-bvc}, \ref{thm:main-civ} \\
\bottomrule
\end{tabular}
\end{center}

\medskip
Falsifiers F5 (epistemic dependence) and F7 (regulatory inadmissibility) are not
directly exercised by the verification ladder because they require external input —
a category-theoretic construction in the case of F5, and a legal ruling in the case
of F7. They are included as honest acknowledgements that the thesis makes claims
(epistemological geometry, regulatory admissibility) whose full falsification space
extends beyond what automated testing can cover.

%%--------------------------------------------------------------------
\section{Openness to Refutation}
\label{sec:openness-refutation}
%%--------------------------------------------------------------------

The Popperian criterion for scientific status is not the absence of falsifiers but the
\emph{genuine possibility of their instantiation}. Each falsifier in this chapter
describes a concrete experiment that a competent engineer could perform. The thesis
does not claim that the falsifiers are impossible — it claims only that they have not
been instantiated as of the 2026-06-09 audit. Any evaluator who instantiates a
falsifier is invited to open an issue in the \wasm{} repository; the authors commit to
investigating within two working days and either refuting the counter-example or
revising the affected theorem.

This commitment is itself falsifiable: if an issue is opened, a counter-example is
demonstrated, and no response is forthcoming within two working days, the thesis
authors have failed their own standard of scientific accountability.
